\subsection{Conventions and normalization}
\label{p:app:conventions}

This appendix collects the conventions of the paper, following~\cite{Erbin:2021:StringFieldTheory}.


\subparagraph{State space.}

The zero modes of the ghosts and of the Virasoro generators are combined into
\begin{equation}
  \label{p:eq:conv:pm-operators}
  L_0^\pm
    := L_0 \pm \bar L_0,
  \qquad
  b_0^\pm
    := b_0 \pm \bar b_0,
  \qquad
  c_0^\pm
    := \frac{1}{2} \, (c_0 \pm \bar c_0).
\end{equation}
$\Hs$ denotes the space of states of the matter--ghost CFT, at every ghost number, that satisfy the level-matching conditions
\begin{equation}
  \label{p:eq:conv:level-matching}
  b_0^- \ket{A}
    = 0,
  \qquad
  L_0^- \ket{A}
    = 0.
\end{equation}
We include all ghost numbers so that the string products below map into $\Hs$.

Ghost numbers are physical: the $\group{SL}(2, \C)$ vacuum has ghost number zero, $c$ counts $+1$ and $b$ counts $-1$.
The tachyon $c_1 \bar c_1 \ket{0}$ then has ghost number two.
The Grassmann parity of a state is its ghost number modulo two:
\begin{equation}
  \label{p:eq:conv:parity}
  \abs{A}
    := \Ngh(A) \bmod 2.
\end{equation}

The $n$-string space $\Hn{n}$ is the space of tensors in $\Hs^{\otimes n}$ whose components are graded-symmetric under permutations of the strands.
Given a basis $\{ \phi_r \}$ of $\mc H$, we can write a state $\Sigma_n \in \mc H^{\wedge n}$ as
\begin{equation}
  \label{p:eq:conv:exp-sigma-n}
  \Sigma_n
    = \sum_{\{ r_i \}} \sigma_{r_1 \cdots r_n} \phi_{r_1} \cdots \phi_{r_n}.
\end{equation}

The parity in which the symmetry is graded depends on the role of the strands.
For the arguments of the vertices and products below, the fields $\Psi^n$ and $\Sigma_n$ and the gauge parameters, the grading is given by the parity $\abs{\phi_{r_i}}$ of the basis state in the strand $i$ of~\eqref{p:eq:conv:exp-sigma-n}, and $\Hn{n}$ is the symmetrized tensor product of $n$ copies of $\mc H$~\cite{Erler:2014:NSNSSectorClosed,Erbin:2020:ClassicalAlgebraicStructures}.
Even strands then have symmetric components and odd strands antisymmetric ones, so that a product is antisymmetric in two gauge parameters.

For the strands that the $n$-string pairing~\eqref{p:eq:conv:pairing-n} sets against a field, those of the outputs of the products and of the multipliers $\Theta_k$, it is the shifted parity $\abs{\phi_{r_i}} + 1$.
This is the dual of the former space under the pairing: if we were also defining the graded-symmetry with respect to the strand parity, the outputs at ghost number three would have antisymmetric components and would pair to zero with $\Sigma_n$.
For this reason, we use instead the shifted parity $\abs{\phi_{r_i}} + 1$.
In this way, the fields have symmetric components, as the pairing requires.
For a tensor with one odd and one even strand, the two parities give the same sign and we need not specify the role.
We write $\Psi^n := \Psi^{\otimes n} \in \Hn{n}$.


\subparagraph{Fields of the paper.}

We list below the fields of the paper with their space, ghost number and Grassmann parity, and the section where each is introduced.
All of them have commuting coefficients, so that the parity of a field is that of its states.
\begin{itemize}
  \item $\Psi \in \Hs$, the classical string field: $\Ngh(\Psi) = 2$, even (\cref{p:sec:hs:action}).

  \item $\Lambda \in \Hs$, the classical gauge parameter: $\Ngh(\Lambda) = 1$, odd (\cref{p:sec:hs:action}).

  \item $\Sigma_k \in \Hn{k} / \ker$, the auxiliary field that replaces the tensor power $\Psi^{\otimes k}$, on the quotient by the kernel of its quadratic operator, $\sprod{2}{2}$ for $k = 2$ and $\wprod{3}{3}$ for $k = 3$: $\Ngh(\Sigma_k) = 2 k$, even (\cref{p:sec:hs:quartic,p:sec:hs:sextic}).
    In the multiplier form of \cref{p:sec:all-orders}, the fields $\Sigma_n \in \Hn{n}$ carry the same letter but live on $\Hn{n}$ itself and are not these auxiliary fields.

  \item $\Theta_k \in \Hn{k}$, the multiplier that enforces $\Sigma_k = \Sigma_{k-1} \otimes \Psi$ in the multiplier form of the action: $\Ngh(\Theta_k) = 3 k$, parity $k \bmod 2$, its strands read in the shifted parity (\cref{p:sec:all-orders}).

  \item $\Upsilon_2 \in \Hn{2} / \ker \sprod{2}{2}$, the partner field: $\Ngh(\Upsilon_2) = 4$, even (\cref{p:app:wrong-sign}).
\end{itemize}
Three letters should not be confused: $d$ is the spacetime dimension, $D$ the ghost dilaton of \cref{p:sec:solutions}, and $\hat d$ the formal dimension of \cref{p:sec:worldsheet}.


\subparagraph{Sphere normalization.}

We normalize the correlation functions on the sphere, with the momentum-conservation factor $(2\pi)^d \delta^{(d)}(\sum k)$ stripped off, as:
\begin{equation}
  \label{p:eq:conv:sphere}
  \braket{c_{-1} \bar c_{-1} c_0^- c_0^+ c_1 \bar c_1}
    = - \frac{1}{2},
  \qquad
  \braket{c_{-1} \bar c_{-1} c_0 \bar c_0 c_1 \bar c_1}
    = - 1.
\end{equation}
The value is real, with no factor $\I$, as the reality of the action under~\eqref{p:eq:conv:reality-fields} requires.
A sphere correlator is nonzero only when the ghost numbers of its insertions add up to six, the ghost-number anomaly.
All the ghost-number formulas below follow from this selection rule.


\subparagraph{Pairing.}

The pairing of two states is the BPZ inner product with one insertion of $c_0^-$:
\begin{equation}
  \label{p:eq:conv:pairing}
  \pair{A}{B}
    := \braket{A | c_0^- | B},
\end{equation}
where $\bra{A}$ is the BPZ conjugate of $\ket{A}$.
We write the BPZ conjugate as $\transp{A}$.
The pairing is nonzero only when $\Ngh(A) + \Ngh(B) = 5$, since the insertion of $c_0^-$ supplies the sixth unit.
Hence it pairs states of opposite parity, $\abs{A} + \abs{B} = 1$.

BPZ conjugation acts on the modes as
\begin{equation}
\label{p:eq:conv:bpz-modes}
\transp{c_n}
  = - c_{-n},
\qquad
\transp{b_n}
  = b_{-n},
\qquad
\transp{(c_0^\pm)}
  = - c_0^\pm,
\end{equation}
and in the same way for $\bar c_n$ and $\bar b_n$.
These are the rules for the inversion $I^+(z) = 1/z$, the standard choice for the closed string~\cite[footnote p.~13]{Zwiebach:1993:ClosedStringField}.
On a monomial of creation modes acting on the vacuum, BPZ conjugation conjugates each mode in place, without a reordering sign:
\begin{equation}
  \label{p:eq:conv:bpz-word}
  \transp{(O_1 \cdots O_k \ket{0})}
    = \bra{0} \transp{O_1} \cdots \transp{O_k}.
\end{equation}

With these rules and the normalization~\eqref{p:eq:conv:sphere}, the pairing is graded-symmetric in the shifted parity:
\begin{equation}
\label{p:eq:conv:pairing-symmetry}
\pair{B}{A}
  = (-1)^{(1 + \abs{A}) (1 + \abs{B})} \, \pair{A}{B}.
\end{equation}
Since the pairing is nonzero only for $\abs{A} + \abs{B} = 1$, the sign is $+1$ on its support: the pairing is symmetric.
This is the rearrangement rule of~\cite{Erbin:2021:StringFieldTheory} restricted to the support.
The sign is checked on explicit states in \cref{p:app:reality}, where the reality conditions are also derived.


\subparagraph{Multi-string pairing.}

On $n$-string states the pairing is the tensor power of~\eqref{p:eq:conv:pairing}, one $c_0^-$ per strand, with the Koszul sign of a tensor product of bilinear forms,
\begin{equation}
  \label{p:eq:conv:pairing-n}
  \pair[n]{A_1 \otimes \cdots \otimes A_n}{B_1 \otimes \cdots \otimes B_n}
    := (-1)^{\sum_{i < j} \abs{B_i} \abs{A_j}} \,
      \pair{A_1}{B_1} \cdots \pair{A_n}{B_n},
\end{equation}
extended bilinearly.
For $n = 1$ it is the pairing~\eqref{p:eq:conv:pairing}, written without index.
When the strands of the first tensor are even, as for $\Psi^n$ and $\Sigma_n$, the sign is $+1$ and the pairing is the plain product of the pairings of the strands.
Each $c_0^-$ thus sits inside the pairing of its own strand, and we read the products in this grouping.
The multi-string pairing is nonzero only when the ghost numbers of the two tensors add up to $5n$.

For every $n$, it is graded-symmetric in the total parity:
\begin{equation}
\label{p:eq:conv:pairing-n-symmetry}
\pair[n]{Y}{X}
  = (-1)^{\abs{X} \abs{Y}} \, \pair[n]{X}{Y},
\qquad
X, Y \in \Hn{n}.
\end{equation}
This follows from combining the Koszul signs with~\eqref{p:eq:conv:pairing-symmetry} strand by strand.
Whenever one of the two tensors is even, as for every pair met in this paper, the sign is $+1$ and the pairing is symmetric.


\subparagraph{String vertices.}

The vertices $\vertex{n} \colon \Hn{n} \to \C$ are the genus-zero string vertices, graded-symmetric multilinear forms.
Their arguments, and those of the products, are graded-symmetric in the parity of the strands for every parity of the arguments.
Hence the product of two gauge parameters in \cref{p:sec:gauge} is antisymmetric in the two.
The outputs of the multi-string products are read in the shifted parity.
The string products $\ell_n \colon \Hn{n} \to \Hs$ are the maps they define through the pairing:
\begin{equation}
  \label{p:eq:conv:vertex}
  \vertex{n}(A_1, \ldots, A_n)
    =: \pair{A_1}{\ell_{n-1}(A_2, \ldots, A_n)},
  \qquad
  \ell_1
    := Q_B.
\end{equation}
For $n = 2$, we have in particular, $\vertex{2}(A_1, A_2) := \pair{A_1}{Q_B A_2}$.
The BRST operator is graded self-adjoint for the pairing:
\begin{equation}
  \label{p:eq:conv:QB-self-adjoint}
  \pair{Q_B A}{B}
    = (-1)^{\abs{A}} \, \pair{A}{Q_B B}.
\end{equation}
This property is the symmetry of $\vertex{2}$.

A vertex extended by multilinearity to tensor arguments is written with the tensors in its slots: using~\eqref{p:eq:conv:exp-sigma-n},
\begin{equation}
  \vertex{n}(\Sigma_k, \cdots)
    := \sum_{\{ r_i \}} \sigma_{r_1 \cdots r_k} \, \vertex{n}(\phi_{r_1}, \ldots, \phi_{r_k}, \cdots).
\end{equation}

The selection rule fixes the ghost number carried by every map.
For $\ell_n$ and $\sprod{m}{n}$ it is the ghost number of the output minus the sum of those of the inputs.
For $\vertex{n}$ it is the ghost number of its antighost insertions, so that $\vertex{n}(A_1, \ldots, A_n)$ is nonzero only when $\sum_i \Ngh(A_i) = 2n$:
\begin{equation}
  \label{p:eq:conv:ghost-numbers}
  \Ngh(\vertex{n})
    = 6 - 2n,
  \qquad
  \Ngh(\ell_n)
    = 3 - 2n,
  \qquad
  \Ngh(\sprod{m}{n})
    = 3n - 2m.
\end{equation}
The last entry concerns the products with several outputs, $\sprod{m}{n} \colon \Hn{m} \to \Hn{n}$, whose output strands have ghost number three.
They are defined by splitting the slots of a vertex into $n$ outputs and $m$ inputs and reading the result through the $n$-string pairing,
\begin{equation}
  \label{p:eq:conv:sprod}
  \vertex{m+n}(A_1, \ldots, A_n, B_1, \ldots, B_m)
    =: \pair[n]{A_1 \otimes \cdots \otimes A_n}{\sprod{m}{n}(B_1, \ldots, B_m)},
\end{equation}
so that $\sprod{n}{1} = \ell_n$ and, for instance:
\begin{equation*}
  \vertex{4}(A_1, A_2, A_3, A_4)
    = \pair[2]{A_1 \otimes A_2}{\sprod{2}{2}(A_3, A_4)}.
\end{equation*}
Each output raises the ghost number by three and each input lowers it by two, so that $\Ngh(\sprod{2}{2}) = 2$, $\Ngh(\sprod{3}{2}) = 0$ and $\Ngh(\sprod{3}{3}) = 3$.

Since $\vertex{m+n}$ is graded-symmetric, reading the same vertex with inputs and outputs exchanged gives the transposition identity
\begin{equation}
  \label{p:eq:conv:transpose}
  \pair[n]{X}{\sprod{m}{n}(Y)}
    = \pair[m]{\sprod{n}{m}(X)}{Y},
  \qquad
  X \in \Hn{n},
  \quad
  Y \in \Hn{m},
\end{equation}
for even arguments, the case of $\Psi$ and the auxiliary fields (odd arguments bring a sign).
$\sprod{n}{m}$ is thus the transpose of $\sprod{m}{n}$ with respect to the multi-string pairing.
We write compositions with $\circ$.

\subparagraph{Normalization of the action.}

With the vertices normalized as in~\eqref{p:eq:conv:vertex}, the action of closed string field theory is:
\begin{equation}
  \label{p:eq:conv:action}
  S
    = \frac{1}{2} \, \pair{\Psi}{Q_B \Psi}
      + \sum_{n \ge 3} \frac{\gs^{n-2}}{n!} \, \vertex{n}(\Psi^n)
    = \sum_{n \ge 1} \frac{\gs^{n-1}}{(n+1)!} \, \pair{\Psi}{\ell_n(\Psi^n)},
\end{equation}
with $\gs$ the string coupling.
The equation of motion and the gauge transformation read:
\begin{equation}
  \label{p:eq:conv:gauge}
  Q_B \Psi
    + \sum_{n \ge 2} \frac{\gs^{n-1}}{n!} \, \ell_n(\Psi^n)
    = 0,
  \qquad
  \delta_\Lambda \Psi
    = Q_B \Lambda
      + \sum_{n \ge 2} \frac{\gs^{n-1}}{(n-1)!} \, \ell_n(\Psi^{n-1}, \Lambda),
\end{equation}
with $\Lambda \in \Hs$ of ghost number one.

The Euclidean path integral is weighted by $\e^{S}$, that is, $S = -S_E$ with $S_E$ the Euclidean action.
  For the tachyon component $t(k)$ of the string field, $\Psi \supset t(k) \, c_1 \bar c_1 \ket{k}$, with $m_T^2 = -4/\alpha'$, the quadratic term of~\eqref{p:eq:conv:action} is:
  \begin{equation}
  \label{p:eq:conv:tachyon-quadratic}
  S
    \supset - \frac{\alpha'}{8} \int \frac{\dd^d k}{(2\pi)^d} \, (k^2 + m_T^2) \, t(k) \, t(-k).
\end{equation}
  It decreases at large Euclidean momentum, as a damped weight requires.
  This fixes the sign of~\eqref{p:eq:conv:sphere} against the sign of the weight.


\subparagraph{Reality.}


We impose reality with the star conjugation of~\cite{Erbin:2021:StringFieldTheory}, the composition of the BPZ and Hermitian conjugations:
\begin{equation}
  \label{p:eq:conv:star}
  A^\star
    := \adj{(\transp{A})}
    = \transp{(\adj{A})}.
\end{equation}
The BPZ bra of $A^\star$ is thus the Hermitian bra of $A$, and conversely:
  \begin{equation}
  \label{p:eq:rea:star-bras}
  \bra{A^\star}
    = \adj{A},
  \qquad
  \adj{(A^\star)}
    = \bra{A},
\end{equation}
  with $\bra{A} = \transp{A}$.
  The two conjugations commute on every mode, since each maps $c_n$, $b_n$ and $\alpha_n$ to a real multiple of the mode of opposite index.
  On the modes, Hermitian conjugation is $\adj{(c_n)} = c_{-n}$ and $\adj{(b_n)} = b_{-n}$, and it reverses the order of a product of modes.
  On operators the star is the same composition, and it distributes over the action on a state with a Koszul sign:
  \begin{equation}
  \label{p:eq:rea:star-operator}
  O^\star
    := \transp{(\adj{O})},
  \qquad
  (O A)^\star
    = (-1)^{\abs{O} \abs{A}} \, O^\star A^\star.
\end{equation}
  The sign appears because $O^\star$ is written to the left of $A^\star$ while Hermitian conjugation reverses the order.
Star conjugation is antilinear and squares to one on every state, even or odd: on a monomial of modes acting on a momentum state it returns the same monomial, up to a sign, on the opposite momentum.
On a tensor it acts strand by strand, without a reordering sign.
On the string field, the reality condition is that of~\cite[eq.~(3.6)]{Zwiebach:1993:ClosedStringField}:
\begin{equation}
  (\ket{\Psi})^\dagger
    = - \bra{\Psi}.
\end{equation}

The pairing is equivariant up to a minus sign,
\begin{equation}
  \label{p:eq:conv:pairing-equivariance}
  \pair{A^\star}{B^\star}
    = - \, \pair{A}{B}^*,
\end{equation}
for all $A$ and $B$.
Complex conjugation of a sphere correlator and Hermitian conjugation of its insertions define two antilinear functionals on the monomials of modes, which vanish on the same monomials.
  Since the vacuum functional is unique up to a factor, there is a constant $\eta$ such that $\braket{W}^* = \eta \, \braket{\adj{W}}$ for every monomial $W$.
  On $W_0 = c_{-1} \bar c_{-1} c_0 \bar c_0 c_1 \bar c_1$, Hermitian conjugation swaps each pair $(c_n, \bar c_n)$, which amounts to three transpositions of anticommuting modes, so that $\adj{W_0} = - W_0$.
  With $\braket{W_0} = - 1$ from~\eqref{p:eq:conv:sphere}, this gives $\eta = - 1$: the vacuum functional is anti-Hermitian, as in~\cite[eq.~(2.37)]{Zwiebach:1993:ClosedStringField}.
  Writing $A = U \ket{0}$ and $B^\star = V \ket{0}$ with monomials of modes $U$ and $V$, the relations~\eqref{p:eq:rea:star-bras} and $\adj{(c_0^-)} = c_0^-$ give $\pair{A^\star}{B^\star} = \eta \, \pair{B}{A}^*$, which is~\eqref{p:eq:conv:pairing-equivariance} by the symmetry~\eqref{p:eq:conv:pairing-symmetry}.
The BRST operator is Hermitian and odd under BPZ conjugation, so that $Q_B^\star = - Q_B$.
  With~\eqref{p:eq:rea:star-operator} and $\abs{Q_B} = 1$, it commutes with star conjugation up to a sign fixed by the parity of its argument:
\begin{equation}
  \label{p:eq:conv:QB-star}
  (Q_B A)^\star
    = (-1)^{\abs{A} + 1} \, Q_B A^\star.
\end{equation}
The vertices are real in the sense that
\begin{equation}
  \label{p:eq:conv:vertex-reality}
  \vertex{n}(A_1, \ldots, A_n)^*
    = (-1)^{n + \sum_{i < j} \abs{A_i} \abs{A_j}} \, \vertex{n}(A_1^\star, \ldots, A_n^\star).
\end{equation}
On even arguments, the sign is $(-1)^n$.
  This property makes the action~\eqref{p:eq:conv:action} real.
For $n = 2$ it follows from~\eqref{p:eq:conv:pairing-equivariance} and~\eqref{p:eq:conv:QB-star}, since $\vertex{2}(A_1, A_2) = \pair{A_1}{Q_B A_2}$.
For $n \ge 3$, we take it as a standing assumption on the vertices.
It follows from the invariance of the vertex regions and measures under complex conjugation of the moduli together with the hermiticity conditions of~\cite[sec.~5.1 and eq.~(7.98)]{Zwiebach:1993:ClosedStringField}.
For the $\group{SL}(2, \C)$ vertices of~\cite{Erbin:2026:SL2CVerticesClosed} it holds for the quartic vertex, and we assume it for every $n$.

Since star conjugation acts strand by strand and preserves the parity of every strand, the multi-string pairing is equivariant as well, $\pair[n]{X^\star}{Y^\star} = (-1)^n \, \pair[n]{X}{Y}^*$.
  The reality conditions on the fields of the paper are:
\begin{equation}
  \label{p:eq:conv:reality-fields}
  \Psi^\star
    = - \Psi,
  \qquad
  \Sigma_k^\star
    = (-1)^k \, \Sigma_k,
  \qquad
  \Lambda^\star
    = - \Lambda.
\end{equation}
The condition on $\Psi$ is the reality condition of the theory, and the condition on $\Lambda$ is the one under which every gauge transformation~\eqref{p:eq:conv:gauge} preserves it~\cite[eq.~(3.18)]{Zwiebach:1993:ClosedStringField}, the free term by the star of $Q_B$ and the interaction terms by the reality of the vertices (\cref{p:app:reality}).
  The condition on $\Sigma_k$ follows from the reality of $\Psi$ and of the vertices once the auxiliary fields are introduced (\cref{p:app:reality}).


\subparagraph{Batalin--Vilkovisky conventions.}

In the Batalin--Vilkovisky formalism of \cref{p:sec:gauge}, every field $\phi^i$ is paired with an antifield $\phi^+_i$ of opposite Grassmann parity and target-space ghost number $\gh(\phi^+_i) = -\gh(\phi^i) - 1$.
The master action is even, of ghost number zero.
In the standard assignment of~\cite{Zwiebach:1993:ClosedStringField}, the antifield of a component of $\Psi$ at worldsheet ghost number $g$ sits at ghost number $5 - g$, so that the two pair under~\eqref{p:eq:conv:pairing}.
For a $2$-string field the pairing is~\eqref{p:eq:conv:pairing-n}, nonzero when the ghost numbers add up to ten, so the antifield of $\Sigma_2$ sits at ghost number six.
We take the antibracket in the convention of~\cite[ch.~15]{Henneaux:1994:QuantizationGaugeSystems},
\begin{equation}
  \label{p:eq:conv:antibracket}
  (F, G)
    := \frac{\pd^r F}{\pd \phi^i} \, \frac{\pd^l G}{\pd \phi^+_i}
      - \frac{\pd^r F}{\pd \phi^+_i} \, \frac{\pd^l G}{\pd \phi^i},
\end{equation}
with right and left derivatives as indicated.

On the string field, the odd symplectic form is the pairing twisted by the parity of its first argument:
\begin{equation}
  \label{p:eq:conv:odd-symplectic}
  \Omega(A, B)
    := (-1)^{\abs{A}} \, \pair{A}{B},
  \qquad
  \Omega(B, A)
    = - (-1)^{\abs{A} \abs{B}} \, \Omega(A, B).
\end{equation}
It is odd and graded-antisymmetric on every pair of states in its support, since the twist changes sign under the exchange of two arguments of opposite parity, whereas the pairing itself is symmetric.
Twisting by the parity of the second argument gives $-\Omega$.
Its overall sign relative to the antibracket~\eqref{p:eq:conv:antibracket} is a convention fixed where the master action is written.

On the auxiliary fields, the parity $\Ngh \bmod 2$ does not give an odd pairing: $\Sigma_2$ and the $2$-string states it pairs with, at ghost numbers four and six, are both even in that grading.
We take the antifield $\Sigma_2^+$ of \cref{p:sec:gauge} to be a $2$-string state paired with $\Sigma_2$ by the $2$-string pairing, hence of worldsheet ghost number six.
It is Grassmann odd and of target-space ghost number $-1$ by the definition of an antifield.
Its worldsheet ghost number is a label on which nothing depends.
The rule~\eqref{p:eq:conv:parity} applies to the string field, its gauge parameter and their antifields.
On the auxiliary pair the overall sign of $\Omega$ is absorbed in the sign of $\Sigma_2^+$ (\cref{p:sec:gauge}).


\subparagraph{Potential.}


The potential is the action at zero momentum with the sign reversed,
\begin{equation}
  \label{p:eq:conv:potential}
  V
    := - S \big\vert_{k = 0},
\end{equation}
so that the tachyon sits at a maximum, $V = -\tfrac{1}{2} \, t^2$ at quadratic order.
A vertex quoted on the potential side carries the sign opposite to that of $\vertex{n}$.
When we count the negative directions of a Hessian, we say on which side it is taken, since the numbers of positive and negative directions of $\Hess S$ are those of $\Hess V$ exchanged.
